\chapter{Hyperparameter Configurations}
\label{app:hyperparameters}

\section{Location of the Hyperparameter Tables}
\label{app:hyper:pointers}

The complete hyperparameter configurations for all models are provided
in the implementation chapter (Chapter~\ref{ch:implementation}):

\begin{itemize}
    \item \textbf{FuzzyArcLoss~V2 pattern classifier:}
    Training configuration including Optuna-optimised scale ($s = 46.1$),
    margin ($m = 0.518$), and threshold ($\tau = 0.448$) is described in
    Section~\ref{sec:impl:pattern} with the full parameter list in the
    accompanying text.

    \item \textbf{ABMIL mutation prediction benchmark:}
    Table~\ref{tab:hyper_abmil2} in
    Section~\ref{sec:impl:mutation:training} lists all hyperparameters
    including learning rate ($2 \times 10^{-4}$), early stopping
    (patience = 15), and maximum tiles per slide (4{,}096).

    \item \textbf{XGBoost baseline (B1):}
    Table~\ref{tab:hyper_xgboost2} in
    Section~\ref{sec:impl:mutation:xgboost} lists
    \texttt{n\_estimators=300}, \texttt{max\_depth=4},
    \texttt{learning\_rate=0.05}.

    \item \textbf{Fuzzy Choquet MIL:}
    Table~\ref{tab:hyper_choquet2} in
    Section~\ref{sec:impl:choquet} lists the 21-parameter
    configuration (6 Shapley + 15 interaction indices).
\end{itemize}

This appendix is intentionally kept as a cross-reference to avoid
duplication with the implementation chapter.


\section{Parameter Count Derivations}
\label{app:param_counts}

The $\sim$200{,}000 and $\sim$266{,}000 parameter counts cited
throughout the thesis are derived below.
All layers use the PyTorch convention:
a fully connected layer from dimension $d_{\text{in}}$ to
$d_{\text{out}}$ has $d_{\text{in}} \times d_{\text{out}} + d_{\text{out}}$
parameters (weights + bias); a layer without bias omits the
$d_{\text{out}}$ term; layer normalisation adds $2 \times d$ parameters
(scale and shift).

\subsection{ABMIL --- Condition P (input\_dim = 518)}

\begin{table}[htbp]
\centering
\small
\caption{Parameter count for Pattern-Informed ABMIL (condition~P,
$\text{input\_dim} = 518$, $\text{hidden} = 256$, $\text{attn} = 128$).}
\label{tab:params_abmil_detail}
\begin{tabular}{l l r}
\hline
\textbf{Component} & \textbf{Computation} & \textbf{Params} \\
\hline
\multicolumn{3}{l}{\textit{Encoder}} \\
Linear($518 \to 256$)       & $518 \times 256 + 256$    & 132{,}864 \\
LayerNorm($256$)            & $2 \times 256$            & 512 \\
\hline
\multicolumn{3}{l}{\textit{Gated Attention}} \\
$\mathbf{V}$: Linear($256 \to 128$) & $256 \times 128 + 128$  & 32{,}896 \\
$\mathbf{U}$: Linear($256 \to 128$) & $256 \times 128 + 128$  & 32{,}896 \\
$\mathbf{w}$: Linear($128 \to 1$, no bias) & $128$             & 128 \\
\hline
\multicolumn{3}{l}{\textit{Classifier}} \\
Linear($256 \to 1$)         & $256 + 1$                 & 257 \\
\hline
\textbf{Total}              &                           & \textbf{199{,}553} \\
\hline
\end{tabular}
\end{table}

For other conditions the encoder input dimension changes:
B2 ($\text{input\_dim} = 512$): 198{,}017 params;
B3 ($\text{input\_dim} = 6$): 68{,}481 params;
A ($\text{input\_dim} = 518$): identical to P.

\subsection{Fuzzy Choquet MIL --- Condition FC (embed\_dim = 512, n\_patterns = 6)}

\begin{table}[htbp]
\centering
\small
\caption{Parameter count for Fuzzy Choquet MIL (condition~FC).
The interpretable fuzzy measure contributes only 42 of the total
$\sim$266{,}000 parameters.}
\label{tab:params_choquet_detail}
\begin{tabular}{l l r}
\hline
\textbf{Component} & \textbf{Computation} & \textbf{Params} \\
\hline
\multicolumn{3}{l}{\textit{Fuzzy Measure (interpretable --- 42 params)}} \\
Shapley values $\boldsymbol{\psi}$       & $6$                        & 6 \\
Interaction matrix $\mathbf{v}_2$        & $6 \times 6$               & 36 \\
\hline
\multicolumn{3}{l}{\textit{Embedding Pathway (shared with ABMIL)}} \\
Encoder: Linear($512 \to 256$)           & $512 \times 256 + 256$     & 131{,}328 \\
Encoder: LayerNorm($256$)                & $2 \times 256$             & 512 \\
$\mathbf{V}_{\text{attn}}$: Linear($256 \to 128$)  & $256 \times 128 + 128$  & 32{,}896 \\
$\mathbf{U}_{\text{attn}}$: Linear($256 \to 128$)  & $256 \times 128 + 128$  & 32{,}896 \\
$\mathbf{w}_{\text{attn}}$: Linear($128 \to 1$, no bias) & $128$            & 128 \\
\hline
\multicolumn{3}{l}{\textit{Merge Projection}} \\
Linear($262 \to 256$)                    & $262 \times 256 + 256$     & 67{,}328 \\
LayerNorm($256$)                         & $2 \times 256$             & 512 \\
\hline
\multicolumn{3}{l}{\textit{Other}} \\
Choquet scale (scalar)                   & $1$                        & 1 \\
Classifier: Linear($256 \to 1$)          & $256 + 1$                 & 257 \\
\hline
\textbf{Total}                           &                           & \textbf{265{,}900} \\
\hline
\end{tabular}
\end{table}

\textbf{Key observation.}
The 42-parameter fuzzy measure (6~Shapley values + 36~interaction
matrix entries, of which 15 upper-triangle entries are used) is the
\emph{only} component that provides intrinsic interpretability
(Shapley values and interaction indices).
The remaining $\sim$265{,}858 parameters constitute the embedding
encoder, gated attention, and merge projection---shared architectural
components that are opaque, as in standard ABMIL.
The merge projection layer ($262 \to 256$, contributing 67{,}840
parameters) accounts for the $\sim$33\% size increase of FC over
ABMIL; it is necessary to fuse the 6-d Choquet output with the
256-d attention-aggregated embedding before classification.
